\chapter{BANACH SPACES}

\section{Natural Transformations}
 Recall that if $V$ is a normed linear space its
 \index{dual!space}%
 \index{space!dual}%
 \index{<unopurstar@$V^*$ (dual space of a normed linear space)}%
\emph{dual space} $V^*$ is the set of all continuous linear functionals on~$V$.  On $V^*$ addition and
scalar multiplication are defined pointwise. The norm is the usual operator norm: if $f \in V^*$, then
 \[ \norm{f} := \inf\{M > 0\colon \abs{f(x)} \le M\norm{x}\quad \text{for all $x \in V$}\} = \sup\{\abs{f(x)}\colon \norm{x} \le 1\}. \]
Under the metric induced by this norm the dual space $V^*$ is complete and thus is a Banach space.

\begin{defn}\label{defn_nls_adjoint} If $T\colon V \sto W$ is a continuous linear map between normed linear spaces,
 \index{<unopurstar@$T^*$ (normed linear space adjoint)}%
$T^*\colon W^* \sto V^*$ is defined as for vector spaces: $T^*(g) = g\,T$ for every
$g \in W^*$.  The map $T^*$, is called the
 \index{adjoint!of a linear map!between normed linear spaces}%
 \index{linear!map!adjoint of a}%
\df{adjoint} of~$T$.
\end{defn}

\begin{exam}\label{exam_dual_functor} The adjoint $T^*$ of a bounded linear map between normed linear spaces $V$ and 
$W$ is itself a bounded linear map between $W^*$ and~$V^*$. In the category of normed linear spaces and continuous linear
 \index{functor!duality}%
 \index{duality functor}%
maps the pair of maps $V \mapsto V^*$ (object map) and $T \mapsto T^*$ (morphism map) is a contravariant
functor from $\cat{NLS_\infty}$ to itself.
\end{exam}

Notice that the preceding definition~\ref{defn_nls_adjoint} differs from the definition of the adjoint of a mapping
between Hilbert spaces (see~\ref{def000926}).  Since every Hilbert space is a normed linear space, this could conceivably
generate confusion.  It is worthwhile convincing oneself of a simple fact: because of the anti-isomorphism (and consequent
isomorphism) between a Hilbert space and its dual provided by the \emph{Riesz-Fr\'echet theorem} (see~\ref{0022057}
and~\ref{cor_Hsp_iso_dual}) the normed linear space adjoint of a bounded linear map is \emph{very} similar (although
certainly not identical) to its Hilbert space adjoint.

\begin{exam}\label{C057717}
The dual $V^{**}$ of the dual $V^*$ of a normed linear space $V$ is the \df{second dual} of~$V$
 \index{dual!second}%
and $T^{**} := \bigl(T^*\bigr)^*$ is a bounded linear map of $V^{**}$ into~$W^{**}$.  On the category
$\cat{BAN_{\boldsymbol\infty}}$ of Banach spaces and bounded linear transformations the
 \index{second!dual!functor for normed linear spaces}%
 \index{functor!second dual!for normed linear spaces}%
\df{second dual functor} comprises the object map $B \mapsto B^{**}$ and the morphism map $T \mapsto T^{**}$.
It is a covariant functor.
\end{exam}

\begin{defn}\label{defn_nat_transf} Let $\cat A$ and $\cat B$ be categories and $F$,$G \colon \cat A \sto \cat B$ be
covariant functors.  A
 \index{natural!transformation}%
 \index{transformation!natural}%
\df{natural transformation} from $F$ to $G$ is a map $\tau$ which assigns to each object
$A$ in $\cat A$ a morphism $\tau_A \colon F(A) \sto G(A)$ in $\cat B$ in such a way that
for every morphism $\alpha \colon A \sto A'$ in $\cat A$ the following diagram commutes.
   \[ \CD
      F(A) @>  F(\alpha)>>  F(A')  \\
         @V\tau_AVV                @VV\tau_{A'}V  \\
      G(A) @>> G(\alpha)>  G(A')
   \endCD \]
We denote such a transformation by $\tau \colon F \sto G$. (The definition of a natural
transformation between two contravariant functors should be obvious: just reverse the
horizontal arrows in the preceding diagram.

A natural transformation $\tau \colon  F \sto G$ is a
 \index{natural!equivalence}%
 \index{equivalence!natural}%
\df{natural equivalence} if each morphism $\tau_A$ is an isomorphism in~$\cat B$.
\end{defn}

\begin{defn}\label{defn_nat_embed} If $V$ is a normed linear space, then the map $\sbsb{\tau}V \colon V \sto V^{**} \colon x
\mapsto x^{**}$ where $x^{**}(f) = f(x)$ for all $f \in V^*$ is called the
 \index{tauv@$\sbsb{\tau}V$ (natural embedding of normed linear spaces)}%
 \index{natural!embedding!for normed linear spaces}%
 \index{embedding!natural}%
\df{natural embedding} of $V$ into~$V^{**}$. (It is not altogether obvious that this mapping is
injective. We will show that it is in proposition~\ref{C067137}.)
\end{defn}

\begin{exam}\label{exam_2dual_nat} In the category $\cat{BAN_{\boldsymbol\infty}}$ of Banach spaces
and bounded linear transformations let $I$ be the identity functor and $(\,\cdot\,)^{**}$ be the
second dual functor.  Then the
 \index{natural!embedding!as a functor}%
 \index{functor!the natural embedding as a}%
 \index{taub@$\tau_B$ (natural embedding of $B$ into $B^{**}$)}%
\emph{natural embedding} $\tau$ is a natural transformation from $I$ to~$(\,\cdot\,)^{**}$.
\end{exam}

\begin{exam} Explain exactly what is meant when people say that the isomorphism
    \[ \fml C(X \uplus Y) \cong \fml C(X) \times \fml C(Y) \]
is a natural equivalence.  The categories involved are $\cat{CpH}$ (compact Hausdorff
spaces and continuous maps) and $\cat{BAN_\infty}$ (Banach spaces and bounded linear
transformations).
\end{exam}

\begin{exam}\label{C057724} On the category $\cat{CpH}$ of compact Hausdorff spaces and
continuous maps let $I$ be the identity functor and let $\fml C^*$ be the functor which
takes a compact Hausdorff space $X$ to the Banach space $\fml (C(X))^*$ and takes a
continuous function $\phi\colon X \sto Y$ between compact Hausdorff spaces to the
contractive linear transformation $\fml C^*\phi = (\fml C\phi)^*$. Define the
 \index{evaluation!map}%
\emph{evaluation map} $E$ by
  \[ E_X\colon X \sto \fml C^*(X)\colon x \mapsto E_X(x) \]
where $(E_X(x))(f) = f(x)$ for all $f \in \fml C(X)$.  (This notation can be cumbersome.  In practice, one usually writes
just $E_x$ for $E_X(x)$.  Thus we have $E_x(f) = f(x)$, which certainly looks better than $(E_X(x))(f) = f(x)$.  In most
contexts it seems unlikely that confusion will result from this simplification of notation.)
 \begin{enumerate}
  \item[(a)]  Notice that this definition does make sense. If $x \in X$, then $E_X(x)$ really
does belong to~$\fml C^*(X)$.
  \item[(b)] Show that the map $E$ makes the following diagram commute
     \[ \xymatrix@+30pt{{X}\ar[r]^*+{\phi} \ar[d]_*+{E_X}
            & {Y} \ar[d]^*+{E_Y} \\
          {\fml C^*(X)} \ar[r]_*+{\fml C^*(\phi)} & {\fml C^*(Y)}   }
     \]
 \end{enumerate}
\end{exam}

The preceding example may seem a bit disappointing.  It seems very much as if we are
attempting to establish a natural transformation between the identity functor (on
$\cat{CpH}$) and the $\fml C^*$ functor.  The problem, of course, is that the two
functors have different target categories.  This is a question we will deal with later.
The evaluation map $E_X\colon X \sto \fml C^*(X)$ is certainly not surjective.
Notice, for example, that every functional of the form $E_X(x)$
preserves multiplication (as well as addition and scalar multiplication). A consequence
of this turns out to be that such functionals all live on the unit sphere of $\fml C^*(X)$.
With a new topology (the so-called \emph{weak-star topology}---see~\ref{C067434}) on
$\fml C^*(X)$ the range of $E_X$ turns out to be a compact Hausdorff space which is naturally
equivalent to $X$ itself.  Even more remarkably, we will eventually show that the range of $E_X$
can be regarded as the maximal ideal space of the algebra~$\fml C(X)$. (See~\ref{000731}.)

\begin{prop}\label{C067137} If $V$ is a normed linear space the natural embedding
$\sbsb{\tau}V \colon V \sto V^{**}\colon x \mapsto x^{**}$ (defined in~\ref{defn_nat_embed})
is an isometric linear map.
\end{prop}

\begin{defn} A Banach space $B$ is
 \index{reflexive!Banach space}%
\df{reflexive} if the natural embedding $\sbsb{\tau}B\colon B \sto B^{**}$ (defined in the preceding proposition) is onto.
\end{defn}

Example~\ref{exam_2dual_nat} shows that in the category of Banach spaces and bounded linear
maps, the mapping $\tau\colon B \mapsto \sbsb{\tau}B$ is a natural transformation from the
identity functor to the second dual functor.  Invoking the \emph{Hahn-Banach theorem} we can
say more.

\begin{prop} In the category of reflexive Banach spaces and bounded linear maps the mapping
$\tau\colon B \mapsto \sbsb{\tau}B$ is a natural equivalence between the identity and second
dual functors.
\end{prop}

\begin{exam} Every Hilbert space is reflexive.
\end{exam}

\begin{exer} Let $V$ and $W$ be normed linear spaces, $U$ be a nonempty convex subset of $V$,
and $f \colon U \sto W$.  Without using any form of the \emph{mean value theorem} show
that if $df = \vc 0$ on $U$, then $f$ is constant on~$U$.  (For the definition of $df$, the
\emph{differential} of~$f$, see sections 13.1--2 in my notes~\cite{Erdman:2007}.)
\end{exer}

\begin{exam}\label{C067181} Let $l_\infty$ be the Banach space of all bounded sequences of
real numbers and $c$ be the subspace of $l_\infty$ comprising all sequences $x = (x_n)$
such that $\lim_{n \sto \infty} x_n$ exists. Define
  \[ f\colon c \sto \R\colon x \mapsto \lim_{n\sto \infty} x_n. \]
Then $f$ is a continuous linear functional on $c$ and $\norm f = 1$.  By the
\emph{Hahn-Banach theorem} $f$ has an extension to $\wh f \in {l_\infty}^*$ with
$\norm{\wh f} = 1$. For every
subset $A$ of $\N$ define a sequence $\bigl(x^A_n\bigr)$ by $x^A_n = \left\{%
\begin{array}{ll}
    1, & \hbox{if $n \in A$;} \\
    0, & \hbox{if $n \notin A$.} \\
\end{array}%
\right.$ Next define $ \mu\colon \sfml P(\N) \sto \R \colon A \mapsto \wh
f\bigl(x^A_n\bigr)$. Then $\mu$ is a finitely additive set function but is not countably
additive.
\end{exam}












\section{Alaoglu's Theorem}
\begin{notn} Let $V$ be a normed linear space, $M \subseteq V$, and $F \subseteq V^\ast$.
Then we define
 \begin{align*}
         M^\perp &:= \{f \in V^*\colon f(x)=0 \text{ for all $x \in M$}\} \\
         F_\perp &:= \{x \in V\colon f(x)=0 \text{ for all $f \in F$}\}
  \end{align*}
One may read $M^\perp$ as ``M perp'' or ``M upper perp''; and $F_\perp$ is ``F perp'' or ``F lower perp''.
The set $M^\perp$ is the
 \index{annihilator}%
 \index{<unopur@$M^\perp$ (annihilator of a set)}%
\df{annihilator} of~$M$ and $F_\perp$ is the
 \index{pre-annihilator}%
 \index{<unopvr@$F_\perp$ (pre-annihilator of a set)}%
\df{pre-annihilator} of~$F$.
\end{notn}

\begin{prop}[Annihilators]  Let $B$ be a Banach space, $M \subseteq B$, and $F \subseteq B^*$. Then
 \begin{enumerate}[label=\rm{(\alph*)}]
  \item If $M \subseteq N \subseteq B$, then $N^\perp \subseteq M^\perp$.
  \item If $F \subseteq G \subseteq B^*$, then $G_\perp \subseteq F_\perp$.
  \item $M^\perp$ is a closed linear subspace of $B^*$.
  \item $F_\perp$ is a closed linear subspace of $B$.
  \item $M \subseteq M^\perp{}_\perp$.
  \item $F \subseteq F_\perp{}^\perp$.
  \item $M^\perp{}_\perp = \bigvee M$.
  \item $M^\perp = B^*$ if and only if $M = \{0\}$.
  \item $F_\perp = B$ if and only if $F = \{0\}$.
  \item If $M$ is a linear subspace of $B$, then $M^\perp = \{0\}$ if and only if $M$ is
dense in $B$.
  \item\label{Fpp_reflexive_case} If $B$ is reflexive, then $F_\perp{}^\perp = \bigvee F$.
 \end{enumerate}
\end{prop}

\begin{proof}[\emph{Hint for proof}] For (k) show that $\tau(F_\perp) = F^\perp$ (where $\tau$ is the
natural embedding) and that therefore $F_\perp{}^\perp = F^\perp{}_\perp$.  \ns
\end{proof}

\begin{exam} Let $B$ be a Banach space that is \emph{not} reflexive.  Let $F = \ran \tau$ (where $\tau$ is the
natural embedding).  Then $F_\perp{}^\perp = B^{**} \ne \bigvee F$, so it is clear that~\ref{Fpp_reflexive_case}
in the preceding proposition need not hold in the nonreflexive case.
\end{exam}

The next result is the Banach space analog of the \emph{fundamental theorem of linear algebra}~\ref{00016025} and
proposition~\ref{prop_ker_adj_perp_ran_op}.

\begin{prop}\label{prop_FTLA_Bsp} Let $T \in \ofml B(B,C)$ where $B$ and $C$ are Banach spaces. Then
 \begin{enumerate}[label=\rm{(\alph*)}]
  \item $\ker T^\ast = (\ran T)^\perp$.
  \item $\ker T = (\ran T^*)_\perp$.
\vskip 3 pt
  \item $\clo{\ran T} = (\ker T^*)_\perp$.
\vskip 3 pt
  \item $\clo{\ran T^*} \subseteq (\ker T)^\perp$.
 \end{enumerate}
\end{prop}

\begin{exam} An example of a Banach space operator $T$ for which $\clo{\ran T^*}$ and $(\ker T)^\perp$ are different
is given in exercise 7, page 243 of~\cite{BrownP:1970}.
\end{exam}

The following example is of course a trivial consequence of proposition~\ref{00015032}, but try to give a simple
direct proof.

\begin{exam}\label{X_l2ball_not_cpt} The closed unit ball in $l_2$ is not compact.
\end{exam}

The usual norm topology on normed linear spaces turns out to be much too large for some applications.  As we have seen
it has so many open sets that the closed unit ball is not compact (except in the finite dimensional case). On the
other hand, too few open sets (the indiscrete topology, as an extreme example) would destroy duality theory by failing
to provide a plentiful supply of continuous linear functionals.  The so-called \emph{weak-star topology} (defined
below) on the dual of a normed linear space turns out to be just the right size.  It has enough open sets to guarantee
a vigorous duality theory and few enough open sets to make the closed unit ball compact
(see \emph{Alaoglu's theorem}~\ref{C067441}).

\begin{defn}\label{C067434} Let $V$ be a normed linear space. The
 \index{weak!topology!on normed spaces}%
 \index{topology!weak!on normed linear spaces}%
\df{weak topology} on $V$ is the weak topology induced by the elements of~$V^*$. (So, of
course, the weak topology on the dual space $V^*$ is the weak topology induced by the
elements of~$V^{**}$.) When a net $(x_\lambda)$ converges weakly to a vector $a$ in $V$
we write
 \index{<arrowright@$x_\lambda \to^w a$ (weak convergence)}%
$x_\lambda \to^w a$. The
 \index{w*@$w^*$-topology (weak star topology)}%
 \index{topology!weak star}%
\df{$w^*$-topology} (pronounced \emph{weak star topology}) on the dual space $V^*$ is the
weak topology induced by the elements of $\ran \tau$ where $\tau$ is the natural
injection of $V$ into~$V^{**}$.  When a net $(f_\lambda)$ converges weakly to a vector $g$
in $V^*$ we write
 \index{<arrowright@$f_\lambda \to^{w^*} g$ (weak star convergence)}%
$f_\lambda \to^{w^*} g$. Notice that when a Banach space $B$ is reflexive (in particular, for Hilbert
spaces) the weak and weak star topologies on $B^*$ are identical.
\end{defn}

\begin{prop} Let $f \in V^*$ where $V$ is a Banach space.  For every $\epsilon > 0$ and every
finite subset $F \subseteq V$ define
  \[ U(f;F;\epsilon)
           = \{g \in V^*\colon \abs{f(x) - g(x)} < \epsilon \text{ for all $x \in F$}\}. \]
The family of all such sets is a base for the $w^*$-topology on~$V^*$.
\end{prop}

The next result, \emph{Alaoglu's theorem}, asserting the $w^*$-compactness of the closed unit ball in the
dual of a normed linear space $V$ is one of the most useful theorems in functional analysis.  Thus
at first it may seem somewhat surprising that it is so hard to find a complete and carefully explained
proof of this result.  The standard proof is in a sense very simple; it is basically just recognizing
that the members of a family $\fml F$ of scalar valued functions can be regarded either
  \begin{enumerate}
     \item[(i)] as functions defined on an uncountable product of compact subsets of the scalar field, or alternatively
     \item[(ii)] as members of the closed unit ball of~$V^*$.
  \end{enumerate}
The crux of the proof is convincing yourself that the convergence of a net of such functions in the product space
is exactly ``the same thing'' as its $w^*$-convergence in the closed unit ball of~$V^*$.  Trying simultaneously to
identify and distinguish between two sets of functions is notationally challenging.  So my best advice is to work through
it yourself until you see how ingenious, but basically simple, the argument is.

\begin{thm}[Alaoglu's theorem]\label{C067441} The closed unit ball of a normed linear space is
 \index{Alaoglu's theorem}%
compact in the $w^*$-topology.
\end{thm}

\begin{proof}[\emph{Hint for proof}] Let $V$ be a normed linear space and denote by $[V^*]_1$ the closed unit ball of its
dual space.  For each $x \in V$ define $D_x = \{\lambda \in \K \colon \abs\lambda \le \norm x\}$.  The product
\smash[b]{$\mathbf P = \prod\limits_{x \in V} D_x$} is compact by \emph{Tychonoff's theorem}.  Consider the function
   \[ \Phi\colon [V^*]_1 \sto \mathbf P \colon f \mapsto \bigl(f(x)\bigr)_{x \in V} \]
where $[V^*]_1$ has the $w^*$-topology and $\mathbf P$ has the product topology (see~\ref{C021437}).  Here we have written
$\bigl(f(x)\bigr)_{x \in V}$ for the indexed family of numbers (or generalized sequence) rather than the more usual
$\bigl(f_x\bigr)_{x \in V}$.  (For more on indexed families see~\cite{Erdman:2007}, section~7.2.)  It is easy to see that
$\Phi$ is injective.  To show that it is continuous let $\bigl(f_\lambda\bigr)_{\lambda \in \Lambda}$ be a net in $[V^*]_1$
and $g$ be a function in $[V^*]_1$ such that $f_\lambda \to^{w^*} g$.  Notice that this happens if and only if
$\pi_x(f_\lambda) \sto \pi_x(g)$ for every $x \in V$ (where $\pi_x$ is a coordinate projection on~$\mathbf P$---see~\ref{C015531}).
Finally, let $\bigl(f_\lambda\bigr)_{\lambda \in \Lambda}$ be a net in $[V^*]_1$ such that $\bigl(\Phi(f_\lambda)\bigr)_{\lambda \in \Lambda}$
converges in~$\mathbf P$.  Show that the net $\bigl(f_\lambda(x)\bigr)_{\lambda \in \Lambda}$ converges in $\K$ for every $x \in V$.
Let $g(x) = \lim\limits_\lambda f_\lambda(x)$ for all $x \in V$ and prove that $g$ is linear, that $\norm g \le 1$, that
$f_\lambda \to^{w^*} g$ in $[V^*]_1$, and that $\Phi(f_\lambda) \sto \Phi(g)$ in~$\mathbf P$.  \ns
\end{proof}

To the best of my knowledge the most reader-friendly version of this standard proof can be found in~\cite{BrownP:1977}, Theorem 15.11.
A much shorter, much easier notationally, much slicker, and much less edifying proof is available in~\cite{Pedersen:1995}, theorem 2.5.2.
I find it less edifying because of its reliance on universal nets.  A net in a set $S$ is
 \index{universal!net}%
 \index{net!universal}%
\df{universal} if for every subset $T$ of $S$ the net is either eventually in $T$ or eventually in its complement. Despite the fact that
every net can be shown to have a universal subnet, even to find a singe nontrivial example of a universal net (that is, one that is not
eventually constant) requires the \emph{axiom of choice}.



















\section{The Open Mapping Theorem}\label{C0694}
A question one can (and should) ask of any concrete category is whether every bijective morphism is necessarily an
isomorphism.  In some categories the answer is trivially \emph{yes} (for example, in $\cat{VEC}$).  In other
categories the answer is trivially \emph{no} (in $\cat{TOP}$, for example, consider the identity map from the
reals with the discrete topology to the reals with their usual topology).  Frequently in categories which arise in
studying analysis, algebra is tugging enthusiastically in one direction while topology is vigorously resisting, pulling
in the other.  Here the question can become quite fascinating---and the answer definitely nontrivial.  In the category
$\cat{BAN_\infty}$ of Banach spaces and continuous linear maps, the question turns out to be intriguingly deep.  In this
case algebra prevails; the answer is affirmative.  It turns out to be a direct consequence of the celebrated \emph{open
mapping theorem}.  Notice as you work through the following proof of this result that for the map in question
completeness of domain and completeness of codomain are both essential---and for quite different reasons.

\begin{defn} A mapping $f \colon X \sto Y$ between topological spaces is
 \index{open!mapping}%
 \index{mapping!open}%
\df{open} if it takes open sets to open sets; that is, $f$ is an open mapping provided that $f^{\sto}(U)$ is open in $Y$
whenever $U$ is open in~$X$.
\end{defn}

The \emph{open mapping theorem} says that continuous linear surjections between Banach spaces are always open mappings.
There is one technical detail that makes the proof of this result a bit tricky.  Suppose that $T$ is a continuous linear
map from a Banach space $B$ into a normed linear space~$V$ and $R$ is the image under $T$ of an open ball around the
origin in~$B$.  We will need to know that if the closure of $R$ contains some open ball about the origin in $V$, then
$R$ itself contains the same ball. It will be convenient to isolate this crucial bit of information in the next proposition.

\begin{notn} In a normed linear space $V$ we denote by
 \index{<unop@$(V)_r$ (open ball of radius $r$ about the origin)}%
$(V)_r$ the open ball about the origin in $V$ of radius~$r$.
\end{notn}

\begin{prop}\label{prop_for_OMT} Let $T\colon B \sto V$ be a continuous linear map from a Banach space into a normed linear space.
If $r$, $s > 0$ and $(V)_s \subseteq \clo{T\bigl(\,(B)_r\,\bigr)}$, then $(V)_s \subseteq T\bigl(\,(B)_r\,\bigr)$.
\end{prop}

\begin{proof}[\emph{Hint for proof}] First verify that without loss of generality we may assume that $r = s = 1$.  Thus
our hypothesis is that $(V)_1 \subseteq \clo{T\bigl(\,(B)_1\,\bigr)}$.

Let $z \in (V)_1$.  The proof is complete if we can show that $z \in T\bigl(\,(B)_1\,\bigr)$.  Choose $\delta > 0$ such
that $\norm z < 1 - \delta$ and let $y = (1 - \delta)^{-1}z$.  Construct a sequence $(v_k)_{k=0}^\infty$ of vectors in~$V$
which satisfy
   \begin{enumerate}
     \item[(i)] $\norm{y - v_k} < \delta^k$ for $k = 0,1,2,\dots$ and
     \item[(ii)] $v_k - v_{k-1} \in T\bigl(\,\delta^{k-1}(B)_1\,\bigr)$ for $k = 1,2,3,\dots$.
   \end{enumerate}
To this end, start with $v_0 = \vc 0$. Use the hypothesis to conclude that $y - v_0 \in \clo{T\bigl(\,(B)_1\,\bigr)}$ and
that therefore there is a vector $w_0 \in T\bigl(\,(B)_1\,\bigr)$ such that $\norm{(y - v_0) - w_o} < \delta$.  Let
$v_1 = v_0 + w_0$ and check that (i) and (ii) hold for $k = 1$.

Then find $w_1 \in T\bigl(\,\delta (B)_1\,\bigr)$ so that $\norm{(y - v_1) - w_1} < \delta^2$ and let $v_2 = v_1 + w_1$.
Check that (i) and (ii) hold for $k= 2$.  Proceed inductively.

For each $n \in \N$ choose a vector $u_n \in \delta^{n-1}(B)_1$ such that $T(u_n) = v_n - v_{n-1}$. Show that the sequence
$(u_n)$ is summable (use\ref{prop_abs_sum_implies_sum}) and let $x = \sum_{n=1}^\infty u_n$.  Conclude the proof by showing
that $Tx = y$  and that $z \in T\bigl(\,(B)_1\,\bigr)$.   \ns
\end{proof}

\begin{thm}[Open mapping theorem]\label{C069414} Every bounded linear surjection between
 \index{open!mapping!theorem}%
 \index{bijective morphisms!are invertible!in $\cat{BAN_1}$}%
Banach spaces is an open mapping.
\end{thm}

\begin{proof}[\emph{Hint for proof}] Suppose that $T\colon B \sto C$ is a bounded linear surjection between Banach spaces.
Start by verifying the following claim: All we need to show is that $T\bigl((B)_1\bigr)$ contains an open ball about the
origin in~$C$.  Suppose for the moment we know that such a ball exists.  Call it~$W$.  Given an open subset $U$ of $B$, we
wish to show that its image under $T$ is open in~$C$.  To this end take an arbitrary point $y$ in $T(U)$ and choose an $x$
in $U$ such that $y = Tx$.  Since $U - x$ is a neighborhood of the origin in $B$, we can find a number $r > 0$ such that
$r\,(B)_1 \subseteq U - x$.  Verify that $T(U)$ is open by proving that $y + rW \subseteq T(U)$.

To establish the claim above, let $V = T\bigl((B)_1\bigr)$ and notice that since $T$ is surjective, $C = T(B) =
\bigcup_{k=1}^\infty k\clo V$.  Use the version of the \emph{Baire category theorem} which says that if a nonempty
complete metric space is the union of a sequence of closed sets, then at least one of these sets has nonempty interior.
(See, for example, Theorem 29.3.21 in\cite{Erdman:2007}.)  Choose $k \in \N$ so that $\bigl(k\clo V\bigr)^o \neq \emptyset$.
Thus there exists $y \in \clo V$ and $r > 0$ such that $y + (C)_r \subseteq k\clo V$.  It is not hard to see from the
convexity of $k\clo V$ that $(C)_r \subseteq k\clo V$.  (Write an arbitrary point $z$ in $(C)_r$ as
$\frac12(y + z) + \frac12(-y + z)$.) Now you have $(C)_{r/k} \subseteq \clo V$.  Use proposition~\ref{prop_for_OMT}.
\ns \end{proof}

\begin{cor}\label{C069417} Every bounded linear bijection between Banach spaces is an
isomorphism.
\end{cor}

\begin{cor} Every bounded linear surjection between Banach spaces
 \index{BAN@$\cat{BAN_\infty}$!quotients in}%
 \index{quotient!map!in $\cat{BAN_\infty}$}%
is a quotient map in~$\cat{BAN_\infty}$.
\end{cor}

\begin{exam} Let $\fml C_u$ be the set of continuous real valued functions on $[0,1]$ under
the uniform norm and $\fml C_1$ be the same set of functions under the $\mathrm{L_1}$ norm
(\,$\norm{f}_1 = \int_0^1\abs{f}\,d\lambda$\,). Then the identity map $I\colon \fml C_u
\sto \fml C_1$ is a continuous linear bijection but is not open. (Why does this not
contradict the \emph{open mapping theorem}?)
\end{exam}

\begin{exam} Let $l$ be the set of sequences of real numbers that are eventually zero,
$V$ be $l$ equipped with the $l_1$ norm, and $W$ be $l$ with the uniform norm.  The
identity map $I\colon V \sto W$ is a bounded linear bijection but is not an isomorphism.
\end{exam}

\begin{exam} The mapping $T\colon \R^2 \sto \R\colon (x,y) \mapsto x$ shows that although
bounded linear surjections between Banach spaces map open sets to open sets they need not
map closed sets to closed sets.
\end{exam}

\begin{prop} Let $V$ be a vector space and suppose that $\norm{\cdot}_1$ and
$\norm{\cdot}_2$ are norms on~$V$ whose corresponding topologies are $\ofml T_1$
and~$\ofml T_2$.  If $V$ is complete with respect to both norms and if $\ofml T_1
\supseteq \ofml T_2$, then $\ofml T_1 = \ofml T_2$.
\end{prop}

\begin{exer} Does there exist a sequence $(a_n)$ of complex numbers satisfying the
following condition: a sequence $(x_n)$ in $\C$ is absolutely summable if and only if
$(a_nx_n)$ is bounded? \emph{Hint.} Consider $ T\colon l_\infty \sto l_1\colon (x_n)
\mapsto (x_n/a_n)$.
\end{exer}

\begin{exam}\label{C069431} Let $\fml C^2([a,b])$ be the family of all twice continuously
differentiable real valued functions on the interval $[a,b]$. Regard it as a vector space
in the usual fashion. For each $f \in \fml C^2([a,b])$ define
    \[ \norm{f} = \norm{f}_u + \norm{f'}_u + \norm{f''}_u\,. \]
This is in fact a norm and under this norm $\fml C^2([a,b])$ is a Banach space.
\end{exam}

\begin{exam}\label{C069432} Let $\fml C^2([a,b])$ be the Banach space given in
exercise~\ref{C069431}, and $p_0$ and $p_1$ be members of $\fml C([a,b])$.   Define
  \[T\colon C^2([a,b]) \sto C([a,b])\colon
                      f \mapsto f'' + p_1f' + p_0f.\]
Then $T$ is a bounded linear map.
\end{exam}

\begin{exer} Consider a differential equation of the form
   \[ y'' + p_1\,y' + p_0\,y = q  \tag{$\ast$} \]
where $p_0$, $p_1$, and $q$ are (fixed) continuous real valued functions on the interval
$[a,b]$.  Make precise, and prove, the assertion that \emph{the solutions to $(\ast)$
depend continuously on the initial values}.  You may use without proof a standard
theorem: For every point $c$ in $[a,b]$ and every pair of real numbers $a_0$ and $a_1$,
there exists a unique solution of $(\ast)$ such that $y(c) = a_0$ and $y'(c) = a_1$.
\emph{Hint.} For a fixed $c \in [a,b]$ consider the map
  \[ S\colon \fml C^2([a,b]) \sto \fml C([a,b]) \times \R^2 \colon
                                        f \mapsto \bigl(Tf,f(c),f'(c)\bigr) \]
where $\fml C^2([a,b])$ is the Banach space of example~\ref{C069431} and $T$ is the
bounded linear map of example~\ref{C069432}.
\end{exer}

\begin{defn} An bounded linear map  $T\colon V \sto W$ between normed linear spaces  is
 \index{bounded!away from zero}%
\df{bounded away from zero} (or
 \index{bounded!below}%
\df{bounded below}) if there exists a number $\delta > 0$ such that $\norm{Tx} \ge
\delta\norm{x}$ for all $x \in V$. (Recall that the word ``bounded'' means one thing when
modifying a linear map and and something quite different when applied to a general
function; so does the expression ``bounded away from zero''. See
definition~\ref{def_bdd_away_zero}.)
\end{defn}

\begin{prop} A bounded linear map between Banach spaces is bounded away from zero if and
only if it has closed range and zero kernel.
\end{prop}




















\section{The Closed Graph Theorem}
A simple necessary and sufficient condition for a linear map $T\colon B \sto C$ between Banach spaces
to be continuous is that its graph be closed.  In one direction this is almost obvious.  The converse is
called the \emph{closed graph theorem}.  When we ask for the graph of $T$ to be
 \index{graph!closed}%
``closed'' we mean closed as a subset of $B \times C$ with the product topology.  Recall (from~\ref{exam_prod_top_norm})
that the topology on the direct sum $B \oplus C$ induced by the product norm is the same as the product topology on $B \times C$.

\begin{exer} If $T\colon B \sto C$ is a continuous linear map between Banach spaces, then its graph is closed.
\end{exer}

\begin{thm}[Closed graph theorem]\label{thm_CGT} Let $T\colon B \sto C$ be a linear map between Banach
 \index{closed!graph theorem}%
spaces.  If the graph of $T$ is closed, then $T$ is continuous.
\end{thm}

\begin{proof}[Hint for proof] Let $G = \{(x,Tx)\colon x \in B\}$ be the graph of~$T$. Apply
(corollary~\ref{C069417} of) the \emph{open mapping theorem} to the map
  \[ \pi\colon G \sto B\colon (x,Tx) \mapsto x\,. \]
\ns \end{proof}

The next proposition gives a very simple necessary and sufficient condition for the graph of a linear map between Banach
spaces to be closed.

\begin{prop}\label{prop_nasc_closed_graph} Let $T\colon B \sto C$ be a linear mapping between Banach spaces. Then $T$ is
continuous if and only if the following condition is satisfied:
   \[ \text{if $x_n \sto 0$ in $B$ and $Tx_n \sto c$ in $C$, then $c = 0$.} \]
\end{prop}

\begin{prop} Suppose $T \colon B \sto C$ is a linear map between Banach spaces such that if
$x_n \sto 0$ in $B$ then $(g \circ T)(x_n) \sto 0$ for every $g \in C^*$.  Then $T$ is
bounded.
\end{prop}

\begin{prop} Let $T \colon B \sto C$ be a linear function between Banach spaces. Define
$T^*f = f \circ T$ for every $f \in C^*$.  If $T^*$ maps $C^*$ into $B^*$, then $T$ is
continuous.
\end{prop}

\begin{defn}\label{C069624} A linear map $T$ from a vector space into itself is
 \index{idempotent}%
\df{idempotent} if $T^2 = T$.
\end{defn}

\begin{prop}\label{prop_cont_idem_op} An idempotent linear map from a Banach space into itself is
continuous if and only if it has closed kernel and range.
\end{prop}

\begin{proof}[\emph{Hint for proof}] For the ``if'' part, use proposition~\ref{prop_nasc_closed_graph}.  \ns
\end{proof}

The next proposition is a useful application of the \emph{closed graph theorem} to Hilbert space operators.

\begin{prop} Let $H$ be a Hilbert space and $S$ and $T$ be functions from $H$ into itself such that
    \[ \langle Sx,y \rangle = \langle x,Ty \rangle \]
for all $x$, $y \in H$.  Then $S$ and $T$ are bounded linear operators.
\end{prop}

\begin{proof}[\emph{Hint for proof}] To verify the linearity of the map $S$ consider the inner product of the vectors
$S(\alpha x + \beta y) - \alpha Sx - \beta Sy$ and $z$, where $x$, $y$, $z \in H$ and $\alpha$, $\beta \in \K$.
To establish the continuity of $S$, show that its graph is closed.  \ns
\end{proof}

Recall that in the preceding proposition the operator $T$ is the \emph{(Hilbert space) adjoint} of $S$;
that is, $T = S^*$.

\begin{defn} Let $a < b$ in~$\R$. A function $f\colon [a,b] \sto \R$ is
 \index{continuously differentiable}%
 \index{differentiable!continuously}%
continuously differentiable if it is differentiable on (an open set containing) $[a,b]$
and its derivative $f\,'$ is continuous on~$[a,b]$.  The set of all continuously
differentiable real valued functions on $[a,b]$ is denoted by $\fml C^1(\,[a,b]\,)$.
\end{defn}

\begin{exam} Let $D\colon \fml C^1(\,[0,1]\,) \sto C(\,[0,1]\,)\colon f \mapsto f\,'$ where
both $C^1(\,[0,1]\,)$ and $C(\,[0,1]\,)$ are equipped with the uniform norm. The mapping
$D$ is linear and has closed graph but is not continuous.  (Why does this not contradict
the \emph{closed graph theorem}?  What can we conclude about~$\fml C^1(\,[0,1]\,)$?)
\end{exam}

\begin{prop} Suppose that $S\colon A \sto B$ and $T\colon B \sto C$ are linear maps between
Banach spaces and that $T$ is bounded and injective.  Then $S$ is bounded if and only if $TS$ is.
\end{prop}

















\section{Banach Space Duality}

\begin{conv}\label{conv_subsp_Bsp} In the context of Banach spaces we adopt the same convention for the use of the word
``subspace'' we used for Hilbert spaces.  The word ``subspace'' will always mean
 \index{subspace!of a Banach space}%
 \index{conventions!subspaces of Banach spaces are closed}%
\emph{closed vector subspace}.  To indicate that $M$ is a subspace of a Banach space $B$ we write
 \index{<binrelle@$\preccurlyeq$ (subspace of a Banach or Hilbert space)}%
$M \preccurlyeq B$.
\end{conv}

\begin{prop} If $M$ is a subspace of a Banach space $B$,
 \index{quotient!Banach space}%
 \index{Banach!space!quotient}%
 \index{space!quotient!Banach}%
then the quotient normed linear space $B/M$ is also a Banach space.
\end{prop}

\begin{proof}[\emph{Hint for proof}] To show that $B/M$ is complete it suffices to show that from an
arbitrary Cauchy sequence $(u_n)$ in $B/M$ we can extract a convergent subsequence.  To this end
notice that if $(u_n)$ is Cauchy, then for every $k \in \N$ there exists $N_k \in \N$ such that
$\norm{u_n - u_m} < 2^{-k}$ whenever $m$, $n \ge N_k$.  Now use induction. Choose $n_1 = N_1$.
Having chosen $n_1$, $n_2$, \dots , $n_k$ so that $n_1 < n_2 < \dots < n_k$ and $N_j \le n_j$ for
$1 \le j \le k$, choose $n_{k+1} = \max\{N_{k+1}, n_k + 1\}$.  This results in a subsequence
$\bigl(u_{n_k}\bigr)$ satisfying $\norm{u_{n_k} - u_{n_{k+1}}} < 2^{-k}$ for all $k \in \N$.
From each equivalence class $u_{n_k}$ choose a vector $x_k$ in $B$ so that $\norm{x_k - x_{k+1}} < 2^{-k+1}$.
Then proposition~\ref{prop_abs_sum_implies_sum} shows that the sequence $(x_k - x_{k+1})$ is
summable.  From this conclude that the sequences $(x_k)$ and $\bigl(u_{n_k}\bigr)$ converge.
(The details of this argument are nicely written out in the proof of theorem 3.4.13 in~\cite{BrownP:1970}.) \ns
\end{proof}

\begin{prop} If $J$ is a closed ideal in a Banach algebra,
 \index{quotient!Banach algebra}%
 \index{Banach!algebra!quotient}%
 \index{algebra!quotient!Banach}%
then $A/J$ is a Banach algebra.
\end{prop}

\begin{thm}[Fundamental quotient theorem for $\cat{BALG}$] Let $A$ and $B$ be Banach algebras
 \index{BALG@$\cat{BALG}$!quotients in}%
 \index{quotient!theorem!for $\cat{BALG}$}%
and $J$ be a proper closed ideal in~$A$. If $\phi$ is a homomorphism from $A$ to $B$ and
$\ker\phi \supseteq J$, then there exists a unique homomorphism $\wt\phi\colon A/J \sto
B$ which makes the following diagram commute.
   \[ \xymatrix@+30pt{A \ar[d]_*+{\pi} \ar[dr]^*+{\phi} & \\
                           A/J \ar[r]_*+{\widetilde\phi} & B  }\]
Furthermore, $\wt\phi$ is injective if and only if $\ker\phi = J$; and $\wt\phi$ is
surjective if and only if $\phi$ is.
\end{thm}

\begin{prop}\label{C037821} Let $J$ be a proper closed ideal in a unital commutative Banach
algebra~$A$.  Then $J$ is maximal if and only if $A/J$ is a field.
\end{prop}

The next proposition specifies quite simple conditions under which a epimorphism between Banach
spaces turns out to be a factor of other morphisms.  At first glance the usefulness of such a
result may seem a bit doubtful.  When you prove proposition~\ref{00151231} notice, however, that
this is \emph{exactly} what you need.

\begin{prop}\label{prop_surj_as_factor} Consider the category of Banach spaces
and continuous linear maps.  If $T\colon A \sto C$, $R\colon A \sto B$, $R$ is
surjective, and $\ker R \subseteq \ker T$, then there exists a unique
$S\colon B \sto  C$ such that $SR = T$. Furthermore, $S$ is injective if and
only if $\ker R = \ker T$.
\end{prop}

\begin{proof}[\emph{Hint for proof}] Consider the following commutative diagram.
Use the \emph{open mapping theorem}.
 \[\xy
    \scalefactor{1.4}%
    \Atrianglepair/>`>`>`<-`>/[A`B`A/\ker R`C;R`\pi`T`\wt R`\wt T]
 \endxy\]   \ns
\end{proof}

It is interesting to note that it is possible to state the preceding result in such a way that
the uniqueness claim is strengthened slightly.  As stated, the conclusion is that there is a
unique bounded linear map $S$ such that $SR = T$.  With no extra work we could have rephrased
the proposition as follows.

\begin{prop}\label{prop_surj_as_factor2} Consider the category of Banach spaces
and continuous linear maps.  If $T\colon A \sto C$, $R\colon A \sto B$, $R$ is
surjective, and $\ker R \subseteq \ker T$, then there exists a unique function
$S$ mapping $B$ into $C$ such that $S\circ R = T$.  The function $S$ must be
bounded and linear.  Furthermore, it is injective if and only if $\ker R = \ker T$.
\end{prop}

Is it worthwhile stating the fussier, second version of the result? It depends, I think, on
where one is headed.  It doesn't seem attractive to opt for the longer version solely on the
grounds that it is more general.  On the other hand, suppose that one later needs the following result.

\begin{prop}\label{prop_factor_blm} If a bounded linear map $T\colon A \sto C$ between Banach spaces can be factored
into the composite of two functions $T = f \circ R$ where $R \in \ofml B(A,B)$ is surjective,
then $f \colon B \sto C$ is bounded and linear.
\end{prop}

\begin{proof}[\emph{Hint for proof}] All that is needed here in order to apply~\ref{prop_surj_as_factor2}
is to notice that $f(0) = 0$ and therefore $\ker R \subseteq \ker T$.   \ns
\end{proof}

Since the preceding is an almost obvious corollary of~\ref{prop_surj_as_factor2}, it is nice to have
this fussier version available.  I suppose that one might argue that while~\ref{prop_factor_blm} may not be
entirely obvious from the statement of~\ref{prop_surj_as_factor}, it is certainly obvious from its
proof, so why bother with the more complicated version?  Such things are surely a matter of taste.

\begin{defn}\label{defn_Bsp_exact_seq}  A sequence of Banach spaces and continuous linear maps
   \[\xymatrix{{\cdots}\ar[r] & B_{n-1}\ar[r]^{T_n}
          & B_n\ar[r]^{T_{n+1}} & B_{n+1}\ar[r] & {\cdots}
   }\]
is said to be
 \index{sequence!exact}%
 \index{exact!sequence}%
\df{exact at} $B_n$ if $\ran T_n = \ker T_{n+1}$. A sequence is \df{exact} if it is
exact at each of its constituent Banach spaces.  A sequence of Banach spaces and
continuous linear maps of the form
   \begin{equation}\label{001500202i}\xymatrix{
      0 \ar[r] & A \ar[r]^S & B\ar[r]^T & C\ar[r] & 0
   }\end{equation}
 \index{sequence!short exact}%
 \index{short exact sequence}%
which is exact at $A$, $B$, and $C$ is a \df{short exact sequence}.  (Here $\vc 0$ denotes the trivial zero dimensional
Banach space, and the unlabeled arrows are the obvious maps.)
\end{defn}

The preceding definitions were for the
 \index{BAN@$\cat{BAN_\infty}$!the category}%
 \index{category!$\cat{BAN_\infty}$ as a}%
category $\cat{BAN_\infty}$ of Banach spaces and continuous linear maps.  Of course the
notion of an \emph{exact sequence} makes sense in many situations---for example, in the
categories of Banach algebras, $C^*$-algebras, Hilbert spaces, vector spaces, Abelian
groups, and modules, among others.

\begin{prop}\label{00151231} Consider the following diagram in the category of Banach spaces
and continuous linear maps
 \[\xymatrix{
      0\ar[r] & A\ar[d]^f\ar[r]^j & B\ar[d]^g\ar[r]^k
              & C\ar@{-->}[d]^h \ar[r] & 0 \\
      0\ar[r] & A'\ar[r]_{j'} & B'\ar[r]_{k'}
      & C'\ar[r] & 0
 }\]
If the rows are exact and the left square commutes, then there exists a unique
continuous linear map $h\colon C \sto C'$ which makes the right square commute.
\end{prop}

\begin{prop}[The Five Lemma]\label{00151232}
 \index{five lemma}%
Suppose that in the following diagram of Banach spaces and continuous linear maps
 \[\xymatrix{
      0\ar[r] & A\ar[d]^f\ar[r]^j & B\ar[d]^g\ar[r]^k
              & C\ar[d]^h \ar[r] & 0 \\
      0\ar[r] & A'\ar[r]_{j'} & B'\ar[r]_{k'}
      & C'\ar[r] & 0
 }\]
the rows are exact and the squares commute.  Then the following hold.
 \begin{enumerate}[label=\rm{(\alph*)}]
  \item If $g$ is surjective, so is $h$.
  \item If $f$ is surjective and $\ran j \supseteq \ker g$, then $h$ is injective.
  \item If $f$ and $h$ are surjective, so is $g$.
  \item If $f$ and $h$ are injective, so is $g$.
 \end{enumerate}
\end{prop}

There are many advantages to working with Hilbert or Banach space operators having closed range.
One immediate advantage is that we recapture the simplicity of the \emph{fundamental theorem of linear algebra}
(\ref{00016025}) from the somewhat more complex versions given in proposition~\ref{prop_ker_adj_perp_ran_op}
and theorem~\ref{prop_FTLA_Bsp}.  On the other hand there are disadvantages.  They do not form a category; the composite
of two closed range operators may fail to have closed range.

\begin{prop} If a bounded linear map $T\colon B \sto C$ between two Banach spaces has closed range, the so does its
adjoint~$T^*$ and $\ran T^* = (\ker T)^\perp$.
\end{prop}

\begin{proof}[\emph{Hint for proof}] Conclude from proposition~\ref{prop_FTLA_Bsp}(d) that one need only show that
$(\ker T)^\perp \subseteq \ran T^*$.  Let $f$ be an arbitrary functional in $(\ker T)^\perp$. What must be found is a
functional $g$ in $C^*$ so that $f = T^*g$.  To this end use the \emph{open mapping theorem} to show that the map $\wt T_\ast$
in the following diagram is invertible. (The maps $T_\ast\colon B \sto \ran T$ and $T$ are exactly the same, except that their
codomains differ.)

 \[\xy
    \scalefactor{1.4}%
    \Atrianglepair/>`>`>`<-`>/[B`\ran T`B/\ker T`\K;T_\ast`\pi`f`\wt{T_\ast}`\wt f]
 \endxy\]

Let $g_0 := \wt f\,{\wt T_\ast}^{-1}$. Use the \emph{Hahn-Banach theorem} to produce the desired~$g$.   \ns
\end{proof}

\begin{prop} In the category $\cat{BAN_\infty}$ if the sequence $\xymatrix{A \ar[r]^S & B \ar[r]^T & C}$ is exact at $B$
and $T$ has closed range, then $\xymatrix{C^* \ar[r]^{T^*} & B^* \ar[r]^{S^*} & A^*}$ is exact at~$B^*$.
\end{prop}

\begin{defn} A functor from the category of Banach spaces and continuous linear maps into itself is
 \index{exact!functor}%
 \index{functor!exact}%
\df{exact} if it takes short exact sequences to short exact sequences.
\end{defn}

\begin{exam}\label{exam_dualfunctor_exact} The functor $B \mapsto B^*$, $T \mapsto T^*$ (see~\ref{exam_dual_functor})
from the category of Banach spaces and continuous linear maps into itself is exact.
\end{exam}

\begin{exam} Let $B$ be a Banach space and $\sbsb{\tau}B\colon x \mapsto x^{**}$ be the natural embedding from $B$ into~$B^{**}$
(see definition~\ref{defn_nat_embed}). Then the quotient space $B^{**}/\ran\sbsb{\tau}B$ is a Banach space.
\end{exam}

\begin{notn} We will denote by $\clo B$ the quotient
space $B^{**}/\ran \sbsb{\tau}B$ in the preceding example.
\end{notn}

\begin{prop} If $T\colon B \sto C$ is a continuous linear map between Banach spaces, then there exists a unique continuous linear map
$\clo T\colon \clo B \sto \clo C$ such that $\clo T(\phi + \ran\sbsb{\tau}B) = (T^{**}\phi) + \ran\sbsb{\tau}C$ for all $\phi \in B^{**}$.
\end{prop}

\begin{proof}[\emph{Hint for proof}] Apply~\ref{00151231} to the following diagram:
 \begin{equation}\label{eqn_exactCD_Bbar}\xymatrix{
      0\ar[r] & B\ar[d]^T\ar[r]^{\sbsb{\tau}B} & B^{**}\ar[d]^{T^{**}}\ar[r]^{\sbsb{\pi}B}
              & \clo B\ar@{-->}[d]^{\clo T} \ar[r] & 0 \\
      0\ar[r] & C\ar[r]_{\sbsb{\tau}C} & C^{**}\ar[r]_{\sbsb{\pi}C} & \clo C\ar[r] & 0
 }\end{equation}  \ns
\end{proof}

\begin{exam}\label{exam_bar_functor} The pair of maps $B \mapsto \clo B$, $T \mapsto \clo T$ from the 
category $\cat{BAN_\infty}$ into itself is an
 \index{<unopu@$B \mapsto \clo B$, $T \mapsto \clo T$ (the functor)}%
 \index{functor!from a Banach space $B$ to $\clo B$}%
exact covariant functor.
\end{exam}

\begin{prop} In the category of Banach spaces and continuous linear maps, if
the sequence
 \[\begin{CD}
  0 \longrightarrow A @> S >> B @> T >> C \longrightarrow 0
  \end{CD}\]
is exact, then the following diagram commutes and has exact rows
and columns.
 \[\bfig
    \iiixiii{'7777}[A`A^{**}`\overline A`B`B^{**}`\overline B`C`C^{**}`\overline C;%
``````S`{S^{**}}`{\overline S}`T`{T^{**}}`{\overline T}]
  \efig\]
\end{prop}

\begin{cor} Every subspace of a reflexive Banach space is reflexive.
\end{cor}

\begin{cor} If $M$ is a subspace of a reflexive Banach space~$B$, then $B/M$ is reflexive.
\end{cor}

\begin{cor} Let $M$ be a subspace of a Banach space~$B$.  If both $M$ and $B/M$ are reflexive,
then so is~$B$.
\end{cor}




















\begin{defn} Let $T \in \ofml L(B,C)$, where $B$ and $C$ are Banach spaces.  Define
$\coker T$, the
 \index{cokernel}%
 \index{cokernel@$\coker T$ (the cokernel of $T$)}%
\df{cokernel} of~$T$, to be $C/\ran T$.
\end{defn}



Unmasking the identity of the dual of a Banach space is frequently a bit challenging.  Following are a few
important examples.  First, however, a word of caution.  Notation for objects living in normed linear spaces
of sequences can be treacherous.  In dealing with completeness of such spaces,one deals with sequences of
sequences of complex numbers.  It is a kindness towards the reader (and likely to yourself a month later
when you try to read something you wrote!) to distinguish notationally between sequences of complex numbers
and sequences of sequences. There are may ways of accomplishing this: double subscripts, subscripts and
superscripts, functional notation for complex sequences, \emph{etc.}

Another problematic notation that arises in sequence spaces is $\sum_{k=1}^\infty a_k \vc e^k$, where $\vc e^k$
is the sequence which is $1$ in the $k^{\text{th}}$ coordinate and $0$ everywhere else. (Recall that in
example~\ref{C063149} we used these vectors for the usual (or standard) orthonormal basis for the Hilbert
space~$l_2$.)  There appears to be little confusion that results from regarding $\sum_{k=1}^\infty a_k \vc e^k$
simply as shorthand for the sequence $(a_k) = (a_1,a_2,a_3,\dots)$ of complex numbers.  If, however, one wishes
to treat it literally as an infinite series, then care is needed.  In the spaces $\sbsb c0$ and $l_1$, there is
no problem. It is indeed the limit as $n$ gets large of $\sum_{k=1}^n a_k \vc e^k$.  (Why?)  But in the spaces
$c$ and $l_\infty$, for example, things go horribly wrong.  To see what can happen, consider the constant sequence
$(1,1,1,\dots)$ (which belongs to both $c$ and $l_\infty$). Then
$\lim_{n \sto\infty}\norm{(1,1,1,\dots) - \sum_{k=1}^n \vc e^k}$ fails to approach zero in either case.
In $c$ the limit is $1$ and in $l_1$ it is~$\infty$.

\begin{exam}\label{exam_dual_C0} The normed linear space $\sbsb c0$ (see~\ref{C0231303})
 \index{c@$c_0$!as a Banach space}%
 \index{Banach!space!$c_0$ as a}%
 \index{l@$l_1$!as the dual of $c_0$}%
 \index{dual!of $c_0$}%
is a Banach space and its dual is $l_1$ (see~\ref{C023191}.
\end{exam}

\begin{proof}[\emph{Hint for proof}] Consider the map
$T\colon {\sbsb c0}^* \sto l_1\colon f \mapsto \sum_{k=1}^\infty f(\vc e^k)\vc e^k$.
(Don't forget to show that $T$ really maps into~$l_1$.) Verify that $T$ is an isometric isomorphism.
To prove that $T$ is surjective, start with a vector $b \in l_1$ and consider the function
$f\colon \sbsb c0 \sto \C\colon a \mapsto \sum_{k=1}^\infty a_kb_k$.    \ns
\end{proof}

\begin{cor} The space $l_1$ is a
 \index{l@$l_1$!as a Banach space}%
 \index{Banach!space!$l_1$ as a}%
Banach space.
\end{cor}

It is interesting that, as it turns out, the spaces $\sbsb c0$ and $c$ are not isometrically isomorphic,
even though their dual spaces are.  The computation of the dual of $c$ is technically slightly fussier than
the one for~${\sbsb c0}^{\!*}$.  It may be helpful to consider the \emph{Schauder bases} for these two spaces.

\begin{defn} A sequence $(\vc e^k)$ of vectors in a Banach space is a
 \index{Schauder basis}%
 \index{basis!Schauder}%
\df{Schauder basis} for the space if for every $x \in B$ there exists a unique sequence $(\alpha_k)$ of
scalars such that $x = \sum_{k=1}^\infty \alpha_k \vc e^k$.
\end{defn}

Two comments about the preceding definition: The word ``unique'' is very important; and it is not true that
every Banach space has a Schauder basis.

\begin{exam} An orthonormal basis in a separable Hilbert space is a Schauder basis for the space.  (See
propositions~\ref{prop_unique_onbasis} and~\ref{prop_Schauder_basis_Hsp}.)
\end{exam}

\begin{exam}\label{exam_dual_c0} We choose the
 \index{usual!Schauder basis for~$\sbsb c0$}%
 \index{Schauder!basis!usual (or standard) for~$\sbsb c0$}%
 \index{basis!standard (or usual)}%
\df{standard} (or
 \index{standard!Schauder basis for $\sbsb c0$}%
\df{usual}) \df{basis vectors} for $\sbsb c0$ in the same way as for~$l_2$ (see~\ref{C063149}).
For each $n \in \N$ let $\vc e^n$ be the sequence in $l_2$ whose $n^{\text{th}}$ coordinate is $1$
and all the other coordinates are~$0$. Then $\{\vc e^n\colon n \in \N\}$ is a Schauder basis
for~$\sbsb c0$.
\end{exam}

At first glance the space $c$ of all convergent sequences may seem enormously larger than $\sbsb c0$,
which contains only those sequences which converge to zero.  After all, corresponding to each vector
$x$ in $\sbsb c0$ there are uncountably many distinct vectors in $c$ (just add any nonzero complex
number to each term of the sequence~$x$).  Is it possible for $c$, containing, as it does, uncountably
many copies of the uncountable collection of members of~$\sbsb c0$, to have a (necessarily countable)
Schauder basis?  The answer, upon reflection, is \emph{yes}.  All we need to do is add a single element
to the basis for~$\sbsb c0$ to deal with the translations.

\begin{exam} Let $\vc e^1$, $\vc e^2$, \dots be as in the preceding example~\ref{exam_dual_c0}.  Also
denote by $\vc 1$ the constant sequence $(1,1,1,\dots)$ in~$c$. Then $(\vc 1, \vc e^1, \vc e^2, \dots)$
is a Schauder basis for~$c$.
\end{exam}

\begin{exam}\label{exam_dual_c} The normed linear space $c$ (see~\ref{C0231302})
 \index{c@$c$!as a Banach space}%
 \index{Banach!space!$c$ as a}%
 \index{l@$l_1$!as the dual of $c$}%
 \index{dual!of $c$}%
is a Banach space and its dual is $l_1$ (see~\ref{C023191}).
\end{exam}

\begin{proof}[\emph{Hint for proof}] Observe initially that for every $f \in c^*$ the sequence
$\bigl(f(\vc e^k)\bigr)$ is absolutely summable and that as a consequence $\sum_{k=1}^\infty a_k f(\vc e^k)$
converges for every sequence $a = (a_1,a_2,\dots)$ in~$c$.  Also note that each $a \in c$ can be written as
$a = a_\infty \vc 1 + \sum_{k=1}^\infty (a_k -a_\infty)\vc e^k$ where $a_\infty := \lim_{k \sto \infty}a_k$
and let $\vc 1 := (1,1,1,\dots)$.  Use this to prove that for every $f \in c^*$
   \[ \abs{f(a)} \le \biggl[ \abs{\beta_f} + \sum_{k=1}^\infty \abs{f(\vc e^k)} \biggr] {\norm a}_\infty \]
where $\beta_f := f(\vc 1) - \sum_{k=1}^\infty f(\vc e^k)$.  For $t = (t_0,t_1,t_2,\dots) \in l_1$ and
$a = (a_1,a_2,a_3,\dots) \in c$ define
   \[ S_t(a) := t_0 a_\infty + \sum_{k=1}^\infty t_ka_k \]
and verify that $S\colon t \mapsto S_t$ is a continuous linear map from $l_1$ into~$c^*$.
Next, for $f \in c^*$ define
   \[ Tf := (\beta_f, f(\vc e^1), f(\vc e^2), \dots) \]
and show that $T\colon f \mapsto Tf$ is a linear isometry from $c^*$ into~$l_1$. Now in order to see that $c^*$ and
$l_1$ are isometrically isomorphic all that is needed is to show that $ST = \id{c^*}$ and $TS = \id{l_1}$.  \ns
\end{proof}











\begin{defn} A topological space is
 \index{locally!compact}%
 \index{compact!locally}%
\df{locally compact} if every point of the space has a compact neighborhood.
\end{defn}

\begin{notn} Let $X$ be a locally compact space.  A scalar valued function $f$ on $X$ belongs to $\fml C_0(X)$ if
 \index{C@$\fml C_0(X)$!space of continuous functions which vanish at infinity}%
it is continuous and vanishes at infinity (that is, if for every number $\epsilon > 0$ there exists a compact subset
of $X$ outside of which $\abs{f(x)} < \epsilon$).
\end{notn}

\begin{exam}\label{exam_C0X} If $X$ is a locally compact Hausdorff space, then, under pointwise algebraic operations and the uniform norm,
 \index{Banach!space!$\fml C_0(X)$ as a}%
 \index{C@$\fml C_0(X)$!as a Banach space}%
$\fml C_0(X)$ is a Banach space.
\end{exam}

We recall one of the most important theorems from real analysis.

\begin{thm}[Riesz representation]\label{C057951} Let $X$ be a locally compact Hausdorff
 \index{Riesz!representation theorem}%
space.  If $\phi$ is a bounded linear functional on $\fml C_0(X)$, then there exists a
unique complex regular Borel measure $\mu$ on $X$ such that
  \[ \phi(f) = \int_X f\,d\mu \]
for all $f \in \fml C_0(X)$ and $\norm\mu = \norm\phi$.
\end{thm}

\begin{proof} This theorem and the following corollary are nicely presented in Appendix C of Conway's
functional analysis text~\cite{Conway:1990}. See also~\cite{Folland:1984}, theorem 7.17; \cite{HewittS:1965},
theorems 20.47 and 20.48; \cite{McDonaldW:1999}, theorem 9.16;and \cite{Rudin:1987}, theorem 6.19.  \ns
\end{proof}

\begin{cor}\label{C057952} If $X$ is a locally compact Hausdorff space, then $\bigl(\fml C_0(X)\,\bigr)^*$
 \index{dual!of $\fml C_0(X)$}%
 \index{C@$\fml C_0(X)$!dual of}%
 \index{M@$\textbf{M}(X)$!as the dual of $\fml C_0(X)$}%
is isometrically isomorphic to~$\textbf{M}(X)$,
 \index{M@$\textbf{M}(X)$!regular complex Borel measures on~$X$}%
the Banach space of all regular complex Borel measures on~$X$.
\end{cor}

\begin{cor} If $X$ is a locally compact Hausdorff space, then $\textbf{M}(X)$
 \index{M@$\textbf{M}(X)$!as a Banach space}%
 \index{Banach!space!$\textbf{M}(X)$ as a}%
is a Banach space.
\end{cor}










In~\ref{defn_nls_adjoint} we defined the adjoint of a linear operator between Banach
spaces and in~\ref{def000926} we defined the adjoint of a Hilbert space
operator. In the case of an operator on a Hilbert space (which is also a Banach space)
it is natural to ask what relationship there is, if any, between these two ``adjoints''.
They certainly can not be equal since the former acts between the dual spaces and the latter
between the original spaces.  In the next proposition we make use of the anti-isomorphism
$\psi$ defined in~\ref{000341} (see~\ref{0022057}) to demonstrate that the two adjoints are
``essentially'' the same.

\begin{prop} Let $T$ be an operator on a Hilbert space $H$ and $\psi$ be the
anti-isomorphism defined in~\ref{000341}.  If we denote (temporarily) the Banach space
dual of $T$ by $T'\colon H^* \sto H^*$, then $T' = \psi\, T^* \psi^{-1}$.  That is, the
following diagram commutes.
   \[ \xy
          \Square/{->}`{<-}`{<-}`{->}/[H^*`H^*`H`H;T'`\psi`\psi`T^*]
      \endxy \]
\end{prop}

\begin{defn}\label{0022091} Let $A$ be a subset of a Banach space $B$.  Then the
 \index{annihilator}%
\df{annihilator} of $A$, denoted
 \index{<perp@$A^\perp$ (annihilator of a set)}%
by~$A^\perp$, is $\{f \in B^*\colon f(a) = 0 \text{ for every $a \in A$}\}$.
\end{defn}

It is possible for the conflict in notations between \ref{0001511} and \ref{0022091} to
cause confusion. To see that in a Hilbert space the annihilator of a subspace and the
orthogonal complement of the subspace are ``essentially'' the same thing, use the
isometric anti-isomorphism $\psi$ between $H$ and $H^*$ discussed in~\ref{0022057} to identify them.

\begin{prop} Let $M$ be a subspace of a Hilbert space $H$. If we (temporarily) denote the
annihilator of $M$ by~$M^a$, then $M^a = \psi^{\sto}\bigl(M^\perp\bigr)$.
\end{prop}
















\section{Projections and Complemented Subspaces}
\begin{defn}
 A bounded linear map from a Banach space into itself is a
 \index{projection!in a Banach space}%
 \index{operator!projection}%
\df{projection} if it is idempotent.  If $E$ is such a map, we say that it is a
projection \df{along $\ker E$ onto $\ran E$}.
\end{defn}

Notice that this is most definitely \emph{not} the same thing as a projection in a
Hilbert space.  Hilbert space projections are \emph{orthogonal projections}; that is,
they are self-adjoint as well as idempotent.  The range of a Hilbert space projection
is perpendicular to its kernel.  Since the terms ``orthogonal'', ``self-adjoint'', and
``perpendicular'' make no sense in Banach spaces, we would not expect the two notions
of ``projection'' to be identical.  In both cases, however, the projections \emph{are}
bounded, linear, and idempotent.

\begin{defn}  Let $M$ be a closed linear subspace of a Banach space~$B$.  If there exists a
closed linear subspace $N$ of $B$ such that $B = M \oplus N$, then we say that $M$ is a
 \index{complemented subspace}%
 \index{subspace!complemented}%
\df{complemented} subspace, that $N$ is its
 \index{complement!of a Banach subspace}%
\df{(Banach space) complement}, and that the subspaces $M$ and $N$ are
\df{complementary}.
\end{defn}

\begin{prop}\label{C069814} Let $M$ and $N$ be complementary subspaces of a Banach space~$B$. Clearly, each
element in $B$ can be written uniquely in the form $m + n$ for some $m \in M$ and $n \in
N$. Define a mapping $E \colon B \sto B$ by $E(m + n) = m$ for $m \in M$ and $n \in N$.
Then $E$ is the projection along $N$ onto~$M$.
\end{prop}

\begin{proof}[\emph{Hint for proof}] Use~\ref{prop_cont_idem_op}.  \ns
\end{proof}

\begin{prop}\label{C069817} If $E$ is a projection on a Banach space $B$, then its kernel and range are
complementary subspaces of~$B$.
\end{prop}

The following is an obvious corollary of the two preceding propositions

\begin{cor}\label{cor_compl_ran_proj} A subspace of a Banach space is complemented if and only if it is
the range of a projection.
\end{cor}

\begin{prop} Let $B$ be a Banach space.  If $E \colon B \sto B$ is a projection onto a subspace
$M$ along a subspace $N$, then $I - E$ is a projection onto $N$ along~$M$.
\end{prop}

\begin{prop}\label{C069824} If $M$ and $N$ are complementary subspaces of a Banach space $B$, then $B/M$ and
$N$ are isomorphic.
\end{prop}

\begin{prop}\label{prop_linv_Bsp} Let $T \colon B \sto C$ be a bounded linear map between Banach spaces.
 \index{left!inverse!of a Banach space operator}%
If $T$ is injective and $\ran T$ is complemented, then $T$ has a left inverse.
\end{prop}

\begin{proof}[\emph{Hint for proof}] For this hint we introduce a (somewhat fussy) notation.
If $f\colon X \sto Y$ is a function between two sets, we define
$f_\bullet \colon X \sto \ran f\colon x \mapsto f(x)$.  That is, $f_\bullet$ is exactly the same
function as $f$, except that its codomain is the range of~$f$.

Since $\ran T$ is complemented there exists a projection $E$ from $C$ onto~$\ran T$
(by~\ref{cor_compl_ran_proj}).  Show that $T_\bullet$ has a continuous inverse.
Let $S = {T_\bullet}^{-1}E_\bullet$.    \ns
\end{proof}

\begin{prop}\label{prop_rinv_Bsp} Let $S \colon C \sto B$ be a bounded linear map between Banach spaces.
 \index{right!inverse!of a Banach space operator}%
If $S$ is surjective and $\ker S$ is complemented, then $S$ has a right inverse.
\end{prop}

\begin{proof}[\emph{Hint for proof}] Since $\ker S$ is complemented there exists $N \preceq C$ such that
$C = N \oplus \ker S$. Show that the restriction of $S$ to $N$ has a continuous inverse and compose this
inverse with the inclusion mapping of $N$ into~$C$.   \ns
\end{proof}

The next result serves as a converse to both propositions~\ref{prop_linv_Bsp} and~\ref{prop_rinv_Bsp}.

\begin{prop}\label{C069831} Let $T \colon B \sto C$ and $S\colon C \sto B$ be bounded linear maps between Banach
spaces.  If $ST = \id B$, then
 \begin{enumerate}[label=\rm{(\alph*)}]
  \item[(a)] $S$ is surjective;
  \item[(b)] $T$ is injective;
  \item[(c)] $TS$ is a projection along $\ker S$ onto~$\ran T$; and
  \item[(d)] $C = \ran T \oplus \ker S$.
 \end{enumerate}
\end{prop}

\begin{prop} If in the category $\cat{BAN_\infty}$ the sequences
 \[ \begin{CD} 0 @>>> A @>j>> B @>f>> C @>>> 0
    \end{CD}\]
 and
 \[ \begin{CD} 0 @>>> C @>g>> B @>p>> D @>>> 0
    \end{CD}\]
are exact and if $fg = \id{C}$, then $pj$ is an isomorphism from $A$ onto~$D$.
\end{prop}

\begin{proof}[\emph{Hint for proof}] Use \ref{C069831} to see that $B = \ker p \oplus \ran j$. Consider the map
$\bigl.p\bigr|_{\ran j}\,j_\bullet$ (notation for $j_\bullet$ as in the hint for~\ref{prop_linv_Bsp}). \ns
\end{proof}

\begin{prop}\label{prop_merge_exactCD} If in the category $\cat{BAN_\infty}$ the diagrams
 \begin{equation} \begin{CD} 0 @>>> A @>j>> B @>f>> C @>>> 0  \\
      @.   @VSVV   @VVUV   @.    @.    \\
 0 @>>> A' @>>j'> B' @>>f'> C' @>>> 0
    \end{CD}\end{equation}
and
 \[ \begin{CD} 0 @>>> C @>g>> B @>p>> D @>>> 0   \\
      @.     @.   @VUVV   @VVTV     @.  \\
  0 @>>> C' @>>g'> B' @>>p'> D' @>>> 0
    \end{CD}\]
commute, if the rows are exact, if $fg = \id{C}$ and if $f'g' = \id{C'}$, then the diagram
 \[ \begin{CD}  A  @>pj>>  D   \\
       @VSVV   @VVTV  \\
     A'  @>>p'j'> D'
    \end{CD}\]
commutes and the maps $pj$ and $p'j'$ are isomorphisms.
\end{prop}

In~\ref{exam_dual_functor} and~\ref{exam_dualfunctor_exact} we saw that in the category of Banach spaces
and continuous maps the duality functor $B \mapsto B*$, $T \mapsto T^*$ is an exact contravariant functor.
And in~\ref{exam_bar_functor} we saw that in the same category the pair of maps $B \mapsto \clo B$,
$T \mapsto \clo T$ is an exact covariant functor.  Composing these two functors creates two new exact 
contravariant functors, which we will denote by $F$ and $G$:
  \[ \xymatrix{B\ar[r]^F &{\clo B}^{\,*}}\,,\,\xymatrix{T\ar[r]^F &{\clo T}^{\,*}} \qquad \text{and} \qquad
           \xymatrix{B\ar[r]^G &\clo{B^*}}\,,\,\xymatrix{T\ar[r]^G &\clo{T^*}}  \]
It is natural to wonder about the relationship between ${\clo B}^{\,*}$ and $\clo{B^*}$.  The next proposition
tells us not only that these two spaces are always isomorphic, but they are naturally isomorphic.  That is, the
 \index{natural!equivalence!of ${\clo B}^{\,*}$ and $\clo{B^*}$}%
functors $F$ and $G$ are naturally equivalent.

\begin{prop} If $T \in \ofml B(B,C)$, where $B$ and $C$ are Banach spaces, then there exist isomorphisms 
$\sbsb{\phi}B \colon {\clo B}^{\,*} \sto \clo{B^*}$ and $\sbsb{\phi}C \colon {\clo C}^{\,*} \sto \clo{C^*}$ such that
   \[ {\clo T}^{\,*} = {\sbsb{\phi}B}^{\!\!\!-1}\,\, \clo{T^*}\,\,\sbsb{\phi}C\,. \]  
\end{prop}

\begin{proof}[\emph{Hint for proof}] Since the duality functor is exact (see~\ref{exam_dualfunctor_exact}) we may 
take the dual of the commutative diagram~\eqref{eqn_exactCD_Bbar} to obtain a second commutative diagram with exact
rows.  Now write down a third commutative diagram by replacing each occurrence of $B$ in the diagram~\eqref{eqn_exactCD_Bbar}
by~$B^*$.  Check that ${\sbsb{\tau}B}^{\!\!*}\,\sbsb{\tau}{B^*} = \id{B^*}$ so that you can use proposition~\ref{prop_merge_exactCD}
on the last two diagrams thus obtaining the commutative diagram
     \[ \xy
          \Square/{<-}`{->}`{->}`{<-}/[{\clo B}^{\,*}`{\clo C}^{\,*}`\clo{B^*}`\clo{C^*};{\clo T}^{\,*}`\sbsb{\phi}B`\sbsb{\phi}C`\clo{T^*}]
      \endxy \]
where $\sbsb{\phi}B := \sbsb{\pi}{B^*}\,{\sbsb{\pi}B}^{\!\!*}$ is an isomorphism.   \ns 
\end{proof}

\begin{cor} A Banach space is reflexive if and only if its dual is.  
\end{cor}

\begin{prop}\label{prop_fd_subsp_complm} Every finite dimensional subspace of a Banach space is complemented.
\end{prop}

\begin{proof}[\emph{Hint for proof}] Let $F$ be a finite dimensional subspace of a Banach space $B$ and let
$\{\vc e^1, \dots, \vc e^n\}$ be a Hamel basis for~$F$.  Then for each $x \in F$ there exist unique scalars
$\alpha_1(x)$, \dots , $\alpha_n(x)$ such that $x = \sum_{k=1}^n \alpha_k(x)\vc e^k$. For each $1 \le k \le n$
verify that $\alpha_k$ is a bounded linear functional, which can be extended to all of~$B$.  Let $N$ be the
intersection of the kernels of these extensions and show that $B = F \oplus N$.  \ns
\end{proof}

\begin{defn} Recall that in~\ref{000082} we defined the \emph{codimension} of a subspace of a vector space
as the dimension of its vector space complement (which always exists by proposition~\ref{0000824}).  We adopt
the same language for Banach spaces: the
 \index{codimension}%
\df{codimension} of any \emph{vector subspace} of a Banach space is the dimension of its \emph{vector space
complement}.
\end{defn}

For the following corollary (of proposition~\ref{prop_fd_subsp_complm}) keep in mind the
convention~\ref{conv_subsp_Bsp} on the use of the word ``subspace'' in the context of Banach spaces.

\begin{cor} In a Banach space every subspace of finite codimension is complemented.
\end{cor}




















Let $M$ be a nonzero subspace of a Banach space~$M$.  Consider the following two properties of the pair $B$ and~~$M$.

\textbf{Property P:} \emph{There exists a projection of norm $1$ from $B$ onto $M$.}

\textbf{Property E:} \emph{For every Banach space $C$, every bounded linear map from $M$ to $C$ can be extended to a bounded
linear map from $B$ to $C$ without increasing its norm.}

\begin{prop} For a nonzero subspace $M$ of a Banach space $B$, properties \textbf{P} and \textbf{E} are equivalent.
\end{prop}

\begin{proof}[\emph{Hint for proof}] To show that \textbf{E} implies \textbf{P} start by extending the identity map
on $M$ to all of~$B$.   \ns
\end{proof}

















\section{The Principle of Uniform Boundedness}
\begin{defn} Let $S$ be a set and $V$ be a normed linear space. A family $\fml F$ of
functions from $S$ into $V$ is
 \index{pointwise!bounded}%
 \index{bounded!pointwise}%
\df{pointwise bounded} if for every $x \in S$ there exists a constant $M_x > 0$ such that
$\norm{f(x)} \le M_x$ for every $f \in \fml F$.
\end{defn}

\begin{defn} Let $V$ and $W$ be normed linear spaces. A family $\fml T$ of bounded linear maps
from $V$ into $W$ is
 \index{uniformly!bounded}%
 \index{bounded!uniformly}%
\df{uniformly bounded} if there exists $M > 0$ such that $\norm T \le M$ for every $T \in
\fml T$.
\end{defn}

\begin{exer}\label{C070114} Let $V$ and $W$ be normed linear spaces. If a family $\fml T$ of
bounded linear maps from $V$ into $W$ is uniformly bounded, then it is pointwise bounded.
\end{exer}

The \emph{principle of uniform boundedness} is the assertion that the converse
of~\ref{C070114} holds whenever $V$ is complete.  To prove this we will make use of a
``local'' analog of this for continuous functions on a complete metric space.  It says, roughly, that if
a family of real valued functions on a complete metric space is pointwise bounded then it is uniformly
bounded on some nonempty open subset of the space.

\begin{prop}\label{C070117} Let $M$ be a complete metric space and $\fml F \subseteq \fml C(M,\R)$.
If for every $x \in M$ there exists a constant $N_x > 0$ such that $\abs{f(x)} \le N_x$ for
every $f \in \fml F$, then there exists a nonempty open set $U$ in $M$ and  a number $N > 0$ such
that $\abs{f(u)} \le N$ for every $f \in \fml F$ and every $u \in U$.
\end{prop}

\begin{proof}[\emph{Hint for proof}] For each natural number $k$ let $A_k =
\bigcap\{f^{\gets}\bigl(\,[-k,k]\,\bigr)\colon f \in \fml F\}$.  Use the same version of the
\emph{Baire category theorem} we used in proving the \emph{open mapping theorem} (\ref{C069414})
to conclude that for at least one $n \in \N$ the set $A_n$ has nonempty
interior. Let $U = \intr{A_n}$.   \ns
\end{proof}

\begin{thm}[Principle of uniform boundedness]\label{thm_PUB} Let $B$ be a Banach space and $W$ be a normed linear space.
 \index{principle!of uniform boundedness}%
 \index{boundedness!principle of uniform}%
 \index{uniform!boundedness!principle of}%
If a family $\ofml T$ of bounded linear maps from $B$ into $W$ is pointwise
bounded, then it is uniformly bounded.
\end{thm}

\begin{proof}[\emph{Hint for proof}] For every $T \in \ofml T$ define
  \[ \sbsb fT\colon B \sto \R\colon x \mapsto \norm{Tx}. \]
Use the preceding proposition to show that there exist a nonempty open subset $U$ of $B$
and a constant $M > 0$ such that $\sbsb fT(x) \le M$ for every $T \in \ofml T$ and every
$x \in U$. Choose a point $a \in B$ and a number $r > 0$ so that $\clo{B_r(a)} \subseteq
U$. Then verify that
  \[ \norm{Ty} \le r^{-1}\bigl(\, \sbsb fT(a + ry) + \sbsb fT(a)\,\bigr) \le 2Mr^{-1} \]
for every $T \in \ofml T$ and every $y \in B$ such that $\norm y \le 1$. \ns
\end{proof}

\begin{thm}[Banach-Steinhaus]  Let $B$ be a Banach space, $W$ be a normed linear space,
 \index{Banach-Steinhaus theorem}%
and $(T_n)$ be a sequence of bounded linear maps from $B$ into~$W$. If the pointwise limit
$\lim_{n \sto \infty} T_nx$ exists for every $x$ in~$B$, then the map $S\colon B \sto
W\colon x \mapsto \lim_{n \sto \infty} T_nx$ is a bounded linear transformation.
\end{thm}

\begin{proof}[\emph{Hint for proof}] For every $x \in B$ the sequence $\bigl(\norm{T_nx}\bigr)_{n=1}^\infty$
is convergent, therefore bounded. Use~\ref{thm_PUB}.   \ns
\end{proof}

\begin{defn} A subset $A$ of a normed linear space $V$ is
 \index{bounded!weakly!set in a normed linear space}%
 \index{weak!boundedness!in a normed linear space}%
 \index{w@$w$-bounded}%
\df{weakly bounded} ($w$-bounded) if $f^{\sto}(A)$ is a bounded subset of the scalar
field of $V$ for every $f \in V^*$.
\end{defn}

\begin{prop} A subset $A$ of a Hilbert space $H$ is weakly bounded if and only if for every $y \in H$
 \index{weak!boundedness!in a Hilbert space}%
the set $\{\langle a,y \rangle \colon a \in A\}$ is bounded.
\end{prop}

The preceding proposition is more or less obvious.  The next one is an important consequence of the
\emph{principle of uniform boundedness}.

\begin{prop}\label{C070131} A subset of a normed linear space is weakly bounded if and only if it
is bounded.
\end{prop}

\begin{proof}[\emph{Hint for proof}] If $A$ is a weakly bounded subset of a normed linear space $V$
and $f \in V^*$, what can you say about $\sup\{a^{**}(f)\colon a \in A\}$?  \ns
\end{proof}

\begin{exer} Your good friend Fred R. Dimm needs help again.  This time he is worried about
the assertion (see~\ref{C070131}) that a subset of a normed linear space is bounded if
and only if it is weakly bounded.  He is considering the sequence $(a^n)$ in the Hilbert
space $l^2$ defined by $a^n = \sqrt{n}\,\vc e^n$ for each $n \in \N$, where $\{\vc
e^n\colon n \in \N\}$ is the usual orthonormal basis for~$l_2$ (see
example~\ref{C063149}).  He sees that it is obviously not bounded (given the usual
topology on~$l^2$).  But it looks to him as if it is weakly bounded.  If not, he argues, then
according to the \emph{Riesz-Fr\'echet theorem} \ref{000342} there would exist a sequence $x$ in
$l^2$ such that the set $\{\abs{\langle a^n,x \rangle}\colon n \in\N\}$ is unbounded. It looks
to him as if this cannot happen because such a sequence would have to decrease more slowly than
the sequence $\bigl(n^{-1/2}\bigr)$, which already decreases too slowly to belong to~$l^2$.
Put Fred's mind to rest by finding him a suitable sequence.  \emph{Hint:} Use a sequence with
lots of zeros.
\end{exer}

\begin{defn} A sequence $(x_n)$ in a normed linear space $V$ is
 \index{weak!Cauchy sequence}%
 \index{Cauchy!sequence!weak}%
 \index{sequence!weak Cauchy}%
\df{weakly Cauchy} if the sequence $(f(x_n))$ is Cauchy for every $f$ in~$V^*$.  The
sequence
 \index{weak!convergence!of sequences in a normed linear space}%
 \index{convergence!weak}%
 \index{sequence!weak convergence of a}%
\df{converges weakly} to a point $a \in V$ if $f(x_n) \sto f(a)$ for every $f \in V^*$
(that is, if it converges in the weak topology induced by the members of~$V^*$). In this
case we write
 \index{<arrowtow@$x_n \to^w a$ (weak sequential convergence)}%
$x_n \to^w a$.  The space $V$ is
 \index{complete!weakly}%
 \index{weak!completeness}%
\df{weakly sequentially complete} if every weakly Cauchy sequence in $V$ converges weakly.
\end{defn}

\begin{prop} In a normed linear space every weakly Cauchy sequence is bounded.
\end{prop}

\begin{prop} Every reflexive Banach space is weakly sequentially complete.
\end{prop}

\begin{exam} Let $f_n(t) = \left\{
                             \begin{array}{ll}
                               1 - nt, & \hbox{if $0 \le t \le \frac1n$;} \\
                               0, & \hbox{if $\frac1n < t \le 1$.}
                             \end{array}
                           \right\}$ for every $n \in \N$. The sequence $(f_n)$ of functions shows that the Banach
space $\fml C([0,1])$ is not w-sequentially complete.
\end{exam}

\begin{cor} The Banach space $\fml C([0,1])$ is not reflexive.
\end{cor}

\begin{prop} Let $\fml E$ be an orthonormal basis for a Hilbert space~$H$ and $(x_n)$ be a
sequence in~$H$. Then $x_n \to^w \vc 0$ if and only if $(x_n)$ is bounded and $\langle
x_n,e \rangle \sto 0$ for every $e \in \fml E$.
\end{prop}

\begin{defn} Let $H$ be a Hilbert space.  A net $(T_\lambda)$ in $\ofml B(H)$
 \index{weak!operator topology!convergence in the}%
 \index{operator!topology!weak}%
 \index{topology!weak operator}%
 \index{weak!convergence!of Hilbert space operators}%
\df{converges weakly} (or \df{converges in the weak operator topology}) to $S \in \ofml B(H)$ if
  \[ \langle T_\lambda x,y \rangle \sto \langle Sx,y \rangle \]
for every $x$, $y \in H$.  In this case we write
$T_\lambda~\to^{\textbf{WOT}}~S$.
 \index{<arrowtowot@$T_\lambda~\to^{\textbf{WOT}}~S$}%

A family $\ofml T \subseteq \ofml B(H)$ is
 \index{weak!boundedness!of a family of Hilbert space operators}%
 \index{bounded!weakly}
 \index{bounded!in the weak operator topology}%
 \index{weak!operator topology!boundedness in}%
 \index{WOT@\textbf{WOT} (weak operator topology)}%
\df{weakly bounded} (or \df{bounded in the weak operator topology}) if for every $x$, $y \in H$
there exists a positive constant $\alpha_{x,y}$ such that
  \[\abs{\langle Tx,y \rangle} \le \alpha_{x,y}\]
for every $T \in \ofml T$.
\end{defn}

It is unfortunate that the preceding definition can be the source of confusion.  Since $\ofml B(H)$ is a
normed linear space it already has a ``weak'' topology (see~\ref{C067434}).  So how does one deal with two
different topologies on $\ofml B(H)$ having the same name?  Some authors call the one just defined the
\emph{weak operator topology} thereby forestalling any possible confusion.  Other writers feel that the
weak topology in the normed linear space sense is of so little use in operator theory on Hilbert spaces
that unless specifically modified the phrase \emph{weak topology} in the context of Hilbert space operators
should be understood to refer to the one just defined. (See, for example, Halmos's comments in the two
paragraphs preceding Problem 108 in~\cite{Halmos:1982}.)  Similarly, often one sees the phrase \emph{a
bounded set of Hilbert space operators.}  Unless otherwise qualified, this is intended to say that the set
in question is \emph{uniformly bounded}.  In any case, be careful of these phrases when reading.

Here is an important consequence of the \emph{principle of uniform boundedness}~\ref{thm_PUB}.

\begin{prop} Every family of Hilbert space operators which is bounded in the weak operator
topology is uniformly bounded.
\end{prop}

\begin{proof}[\emph{Hint for proof}] Let $H$ be a Hilbert space and $\ofml T$ a family of operators on~$H$
which is bounded with respect to the weak operator topology. Use proposition~\ref{C070131} to show, for a
fixed $y \in H$, that $\{T^*y\colon T \in \ofml T\}$ is a bounded subset of~$H$.  Use this (and~\ref{C070131}
again) to show that $\{Tx\colon \text{$T \in \ofml T$, $x \in H$, and $\norm x \le 1$}\}$ is bounded in~$H$. \ns
\end{proof}

As noted above, many writers would shorten the statement of the preceding proposition to:
\emph{Every weakly bounded family of Hilbert space operators is bounded.}

\begin{defn}
A net $(T_\lambda)$ in $\ofml B(H)$
 \index{strong!operator topology!convergence in the}%
 \index{operator!topology!strong}%
 \index{topology!strong operator}%
 \index{strong!convergence!of Hilbert space operators}%
 \index{SOT@\textbf{SOT} (strong operator topology)}%
\df{converges strongly} (or \df{converges in the strong operator topology}) to $S \in \ofml B(H)$ if
  \[ T_\lambda x \sto Sx \]
for every $x \in H$.  In this case we write $T_\lambda~\to^{\textbf{SOT}}~S$.
 \index{<arrowtowot@$T_\lambda~\to^{\textbf{SOT}}~S$}%

The usual norm convergence of operators in $\ofml B(H)$ is often referred to as
 \index{uniform!operator topology!convergence in the}%
 \index{operator!topology!uniform}%
 \index{topology!uniform operator}%
 \index{uniform!convergence!of Hilbert space operators}%
\df{uniform convergence} (or \df{convergence in the uniform operator topology}).  That is, we write
 \index{<arrowtowot@$T_\lambda~\sto~S$}%
$T_\lambda \sto S$ if $\norm{S - T_\lambda}\sto 0$.
\end{defn}

\begin{prop} Let $H$ be a Hilbert space, $(T_n)$ be a sequence in $\ofml B(H)$,
and $B \in \ofml B(H)$. Then the following implications hold:
   \[ T_n \sto B \quad \implies \quad T_n~\to^{\textbf{SOT}}~B
                                \quad \implies \quad T_n~\to^{\textbf{WOT}}~B. \]
\end{prop}

\begin{prop} Let $(S_n)$ and $(T_n)$ be sequences of operators on a Hilbert space~$H$. \\
If $S_n \to^{\textbf{WOT}} A$ and $T_n \to^{\textbf{SOT}} B$, then $S_nT_n \to^{\textbf{WOT}} AB$.
\end{prop}

\begin{proof}[\emph{Hint for proof}] Write $S_nT_n - AB$ as $S_nT_n - S_nB + S_nB - AB$.  \ns
\end{proof}

\begin{exam} In the preceding proposition the hypothesis $T_n \to^{\textbf{SOT}} B$ cannot
be replaced by $T_n~\to^{\textbf{WOT}}~B$.
\end{exam}

\begin{proof}[\emph{Hint for proof}] Consider powers of the unilateral shift operator.  \ns
\end{proof}







\endinput

